% !TeX root = ../MAgent_Theory_2.0.tex
\chapter{Connections in a Moving Presentation}
\label{app:connection}

Links between agents at a single context do not determine transport between neighboring contexts. A connection requires a smooth dependence on the direction and size of a base displacement. This appendix records, in matrix form, the transformation law of connection data supplied in a moving presentation: given a presentation $U_i$ and a one-form $a_i$, it computes the displayed connection and curvature and their behavior under changes of coordinates. The one-form $a_i$ is supplied for the transformation calculation in this appendix. \Cref{ch:link-base-transport} separately constructs it from regular contextual laws, interfaces, and links. The remaining selection and continuation questions are \cref{open:conn-derived}. Collective least-squares selection in \cref{geo:collective} is a separate construction. The context-local and formation flows use it diagnostically. The added energy in \cref{ch:contextual-coupling} instead makes collective alignment enter the law velocity through a contextual force.

Fix one sector and one agent, suppressing the sector superscript. On a coordinate neighborhood of $\mathcal C$, let $U_i(c)$ be a twice continuously differentiable invertible presentation and let
$a_i=a_{i,\mu}(c)dc^\mu$ be a continuously differentiable matrix-valued one-form. For a principal $G$-connection its values must lie in the represented Lie algebra. General endomorphism-valued forms instead describe connections for the extended general linear frame bundle.

\definitionheading{Infinitesimal directional transport}{def:conn-directional}
The proposed transport family has first-order behavior
\begin{equation}
\Theta_i(c+\delta c,c)
=U_i(c+\delta c)\exp[-a_{i,\mu}(c)\delta c^\mu]U_i(c)^{-1}
+O(|\delta c|^2).
\label{eq:conn-infinitesimal}
\end{equation}
The first-order generator is linear in the tangent displacement. A single finite matrix at $c$ supplies neither this linear dependence nor the identity value at zero displacement.

\propositionheading{Connection and curvature}{prop:conn-curvature}
With the convention $D_i=d+B_i$, \cref{eq:conn-infinitesimal} has connection and curvature
\begin{equation}
B_i=U_i a_iU_i^{-1}-dU_iU_i^{-1},\qquad
F_i=U_i(da_i+a_i\wedge a_i)U_i^{-1}.
\label{eq:conn-curvature}
\end{equation}
\begin{proof}
Expanding the two factors at $c$ gives
$\Theta_i=I+(dU_iU_i^{-1}-U_i a_iU_i^{-1})(\delta c)+O(|\delta c|^2)$,
which identifies $B_i$. Acting on a vector-valued function gives the operator identity
$d+B_i=U_i(d+a_i)U_i^{-1}$. Squaring it yields
$(d+B_i)^2=U_i(d+a_i)^2U_i^{-1}$, and the squares act by the curvature two-forms in \cref{eq:conn-curvature}.
\end{proof}

For an explicit curved example, take $U_i=I$ over an open subset of $\R^2$ and set
\begin{equation}
a_{i,x}=E_{12},\qquad a_{i,y}=E_{21}.
\end{equation}
Both are supplied rank-one receiving--sending products, while their derivatives vanish. Nevertheless,
\begin{equation}
F_{i,xy}=[E_{12},E_{21}]
=\begin{pmatrix}1&0\\0&-1\end{pmatrix}\ne0.
\end{equation}
Thus even constant directional interfaces can give curvature through noncommutativity. The example is supplied. Its curvature has the real eigenvalues $\pm1$, and the spectrum of a curvature value is unchanged by every change of frame. A connection that preserves a fiber metric $M$, $dM=Ma+a^TM$, has curvature skew for $M$ and therefore purely imaginary spectrum, so this example is not produced by any selector whose output preserves a fiber metric, such as the rectangular selector of \cref{ch:link-base-transport} with its induced metric \cref{ibt:eq:metric-compatibility}. By contrast, replacing a finite pointwise map by a new frame $V(c)$ and setting $B=-dV V^{-1}$ gives zero curvature: it is the case $a=0$ of \cref{eq:conn-curvature}.

Finite paths use the solution of
\begin{equation}
\frac{d}{dt}T_i(t)=-B_i(\dot\gamma(t))T_i(t),\qquad T_i(0)=I.
\end{equation}
Uniqueness of this linear equation gives composition under concatenation and inversion under path reversal. The solution is the path-ordered exponential; the usual link and Wilson-line constructions use this ordering \citep{Tong2018}. Independently declared opposite attending edges are not two orientations of this same path and need not be inverses.

\propositionheading{Compatibility under re-expression}{prop:conn-patches}
Under an agent-coordinate change $U_i'=g_iU_i$ with the same reference form,
\begin{equation}
B_i'=g_iB_i g_i^{-1}-dg_i g_i^{-1}.
\end{equation}
Under a change of reference presentation $U_i'=U_i h$, which sends reference laws to their $h^{-1}$ pushforwards and is the inverse of a target change written $p'=h\cdot p$, keeping the same intrinsic transport instead requires
\begin{equation}
a_i'=h^{-1}a_i h+h^{-1}dh.
\label{eq:conn-reference-affine}
\end{equation}
\begin{proof}
Insert $g_iU_i$ into \cref{eq:conn-curvature} and apply the product rule. Inserting $U_i h$ and \cref{eq:conn-reference-affine} cancels the two terms containing $dh$, leaving the original $B_i$.
\end{proof}
The same affine law must hold between overlapping coordinate patches. A directional interface product transforming only by conjugation is a tensorial one-form. It can describe the difference between two connections, but it is not by itself the full affine connection data. Constructing displacement links with the endpoint transformation law and then differentiating them is one way to obtain the required inhomogeneous term.

If patch coordinates satisfy $v_\alpha=g_{\alpha\beta}v_\beta$, their transition functions must obey $g_{\alpha\beta}g_{\beta\gamma}=g_{\alpha\gamma}$ on triple overlaps. The displayed affine law then makes the local connections compatible. A nontrivial bundle need not admit a single global presentation $U_i$. Even on a trivial bundle, zero curvature does not imply path independence: the scalar connection $B=b d\theta$ on the circle is flat but has holonomy $e^{-2\pi b}$, which differs from one for real $b\ne0$. Global path independence requires trivial holonomy as well as local flatness.

Different agents may be assigned different connection data in their own presentations. Since all agents are sections of bundles associated to one principal bundle $P$, and in the linear realization of one associated vector bundle, these data define one connection on $P$ only when, after re-expression in a common frame, their local forms agree and satisfy its overlap relations; otherwise they are several connections on the same bundle, each used by one agent. Neither case requires an averaging construction.

\openproblemheading{Selecting and continuing interaction-derived transport}{open:conn-derived}
The full-rank construction in \cref{ibt:displacement,ibt:path} supplies displacement identity, a linear first jet, compatible endpoint laws, and path transport from contextual law and link messages. The regular alignment construction in \cref{ibt:alignment} gives another selector, and its reciprocal-link example has nonzero curvature. Both require a declared comparison rule. The directed backtrack sector can enter the data and is not identified with the reversal of a contextual path.

The remaining problem is to select the message labels, column weights, or augmented comparison metric from a stated agent model, and to establish continuation when rank or infinitesimal identifiability fails. Link-only data cannot identify a continuous gauge subgroup that fixes every input, as \cref{ibt:invisible} proves. A selector also needs new dynamical analysis before it can replace the law-Fisher auxiliary field in \cref{ch:contextual-coupling}: its full constitutive derivative must enter the energy variation. Neither construction derives the contextual base, its topology, or its dimension.
